\subsection{存在与唯一性定理}

设 \(\lambda \in \mathfrak{h}^*\) 。由于 \(\mathfrak{b}_{+} = \mathfrak{h} \oplus \mathfrak{n}_{+}\) 且 \(\mathfrak{n}_{+} = [\mathfrak{b}_{+}, \mathfrak{b}_{+}]\) ，映射 \(h + n \mapsto \lambda(h)\) （其中 \(h \in \mathfrak{h}\) ，\(n \in \mathfrak{n}_{+}\) ）（从 \(\mathfrak{b}_{+}\) 到 \(k\) 的）是 \(\mathfrak{b}_{+}\) 的一个 1 维表示。用 \(\mathrm{L}_{\lambda}\) 表示 \(k\) -向量空间 \(k\) ，它配备了由这个表示定义的 \(\mathfrak{b}_{+}\) -模结构。设 \(\mathrm{U}(\mathfrak{g})\) ，\(\mathrm{U}(\mathfrak{b}_{+})\) 是 \(\mathfrak{g}\) ，\(\mathfrak{b}_{+}\) 的包络代数，于是 \(\mathrm{U}(\mathfrak{b}_{+})\) 是 \(\mathrm{U}(\mathfrak{g})\) 的一个子代数；回顾 \(\mathrm{U}(\mathfrak{g})\) 是一个自由右 \(\mathrm{U}(\mathfrak{b}_{+})\) -模（第一章 §2 第 7 节 定理 1 的推论 5）。置

\[
\mathrm{Z} (\lambda) = \mathrm{U} (\mathfrak {g}) \otimes_ {\mathrm{U} (\mathfrak {b} _ {+})} \mathrm{L} _ {\lambda}.\tag{2}
\]

则 \(Z(\lambda)\) 是一个左 \(\mathfrak{g}\) -模。用 \(e\) 表示元素 \(1 \otimes 1\) （\(Z(\lambda)\) 中的）。

命题 5. (i) 元素 \(e\) （\(Z(\lambda)\) 的）是权为 \(\lambda\) 的本原元，且生成 \(\mathfrak{g}\) -模 \(Z(\lambda)\) 。

\begin{enumerate}
\def\labelenumi{(\roman{enumi})}
\setcounter{enumi}{1}
\item
  设 \(\mathrm{Z}^{+}(\lambda) = \sum_{\lambda \neq \mu} \mathrm{Z}(\lambda)^{\mu}\) 。\(\mathrm{Z}(\lambda)\) 的每个异于 \(\mathrm{Z}(\lambda)\) 的子模都包含在 \(\mathrm{Z}^{+}(\lambda)\) 中。
\item
  存在最大子模 \(F_{\lambda}\) （\(Z(\lambda)\) 的，异于 \(Z(\lambda)\) 的）。商模 \(Z(\lambda)/F_{\lambda}\) 是单的且具有最高权 \(\lambda\) 。
\end{enumerate}

显然 e 生成 g-模 \(Z(\lambda)\) 。若 \(x \in b_{+}\) ，

\[
x. e = (x. 1) \otimes 1 = (1. x) \otimes 1 = 1 \otimes x. 1 = \lambda (x) (1 \otimes 1) = \lambda (x) e,
\]

由此得 (i)。

\(\mathfrak{h}\) -模 \(Z(\lambda)\) 是半单的（命题 1）。若 \(G\) 是一个 \(\mathfrak{g}\) -子模（\(Z(\lambda)\) 的），则 \(G = \sum_{\mu \in \mathfrak{h}^*}(G \cap Z(\lambda)^\mu)\) 。假设 \(G \cap Z(\lambda)^\lambda \neq 0\) 蕴含 \(G = Z(\lambda)\) ，因为 \(\dim Z(\lambda)^\lambda = 1\) 且 \(e\) 生成 \(\mathfrak{g}\) -模 \(Z(\lambda)\) 。若 \(G \neq Z(\lambda)\) ，则 \(G = \sum_{\mu \neq \lambda} G \cap Z(\lambda)^\mu \subset Z^+( \lambda)\) 。

设 \(F_{\lambda}\) 是 \(Z(\lambda)\) 的异于 \(Z(\lambda)\) 的诸 g-子模之和。由 (ii)，\(F_{\lambda} \subset Z^{+}(\lambda)\) 。故 \(F_{\lambda}\) 是 \(Z(\lambda)\) 的异于 \(Z(\lambda)\) 的最大子模。显然 \(Z(\lambda)/F_{\lambda}\) 是单的，且 e 在 \(Z(\lambda)/F_{\lambda}\) 中的典则像是权为 \(\lambda\) 的本原元。

在本章的其余部分，命题 5 的 g-模 \(Z(\lambda)/F_{\lambda}\) 将记作 \(E(\lambda)\) 。

定理 1. 设 \(\lambda \in \mathfrak{h}^*\) 。\(\mathfrak{g}\) -模 \(\mathrm{E}(\lambda)\) 是单的且具有最高权 \(\lambda\) 。每个单 \(\mathfrak{g}\) -模若其最高权为 \(\lambda\) ，便同构于 \(\mathrm{E}(\lambda)\) 。

第一个断言由命题 5 (iii) 得出。第二个由命题 4 得出。

命题 6. 设 V 是一个 g-模，\(\lambda\) 是 \(h^{*}\) 的一个元素，v 是 V 的权为 \(\lambda\) 的本原元。

\begin{enumerate}
\def\labelenumi{(\roman{enumi})}
\item
  存在唯一的 \(\mathfrak{g}\) -模同态 \(\psi: \mathrm{Z}(\lambda) \to \mathrm{V}\) 使 \(\psi(e) = v\) 。
\item
  假设 v 生成 V。则 \(\psi\) 是满射。此外，\(\psi\) 是双射当且仅当对 \(\mathrm{U}(\mathfrak{n}_{-})\) 的每个非零元素 u ，\(u_{V}\) 是单射。
\item
  映射 \(u \mapsto u \otimes 1\) （从 \(\mathrm{U}(\mathfrak{n}_{-})\) 到 \(\mathrm{Z}(\lambda)\) 的）是双射。
\end{enumerate}

设 K 是 \(\mathrm{U}(\mathfrak{b}_{+})\) 在 \(L_{\lambda}\) 上的表示的核；它在 \(\mathrm{U}(\mathfrak{b}_{+})\) 中的余维数是 1 。设 \(\mathrm{J} = \mathrm{U}(\mathfrak{g})\mathrm{K}\) 是由 K 生成的 \(\mathrm{U}(\mathfrak{g})\) 的左理想；于是 \(L_{\lambda}\) 可以与 \(\mathrm{U}(\mathfrak{b}_{+})/\mathrm{K}\) 等同（作为左 \(\mathrm{U}(\mathfrak{b}_{+})\) -模），且 \(\mathrm{Z}(\lambda)\) 可以与 \(\mathrm{U}(\mathfrak{g})/\mathrm{J}\) 等同（作为左 \(\mathrm{U}(\mathfrak{g})\) -模）。我们有 K.v = 0 ，故 J.v = 0 ，这就证明了 (i)。

现在假设 \(v\) 生成 V。显然 \(\psi\) 是满射。

由 Poincaré-Birkhoff-Witt 定理（第一章 §2 第 7 节 定理 1 的推论 6），\(\mathrm{U}(\mathfrak{n}_{-})\) 在 \(k\) 上的一个基也是 \(\mathrm{U}(\mathfrak{g})\) 作为右 \(\mathrm{U}(\mathfrak{b}_{+})\) -模的一个基。故映射 \(\varphi : u \mapsto u \otimes 1\) （从 \(\mathrm{U}(\mathfrak{n}_{-})\) 到 \(\mathrm{U}(\mathfrak{g}) \otimes_{\mathrm{U}(\mathfrak{b}_{+})} \mathrm{L}_{\lambda}\) 的）是双射。设 \(u \in \mathrm{U}(\mathfrak{n}_{-})\) 。则 \(\varphi^{-1} \circ u_{\mathrm{Z}(\lambda)} \circ \varphi\) 是由 \(u\) 在 \(\mathrm{U}(\mathfrak{n}_{-})\) 上作的左乘。鉴于第一章 §2 第 7 节 定理 1 的推论 7，\(u_{\mathrm{Z}(\lambda)}\) 是单射若 \(u \neq 0\) 。于是，若 \(\psi\) 是双射，则 \(u_{\mathrm{V}}\) 是单射，只要 \(u\) 是 \(\mathrm{U}(\mathfrak{n}_{-})\) 中的非零元。

假设 \(\psi\) 不是单射。存在 \(u \in \mathrm{U}(\mathfrak{n}_{-})\) 使 \(u \neq 0\) 且 \(\psi(\varphi(u)) = 0\) 。则

\[
u _ {\mathrm{V}}. v = u _ {\mathrm{V}}. \psi (1 \otimes 1) = \psi (u \otimes 1) = \psi (\varphi (u)) = 0.
\]

推论 1. 设 \(\lambda \in \mathfrak{h}^*\) 与 \(\alpha \in B\) 使 \(\lambda(H_{\alpha}) + 1 \in \mathbf{N}\) 成立。则 \(Z(-\alpha + s_{\alpha}\lambda)\) 同构于一个 \(\mathfrak{g}\) -子模（\(Z(\lambda)\) 的）。

置 \(m = \lambda(H_{\alpha})\) 。设 \(x = X_{-\alpha}^{m+1}.e \in Z(\lambda)\) ，并设 V 是由 x 生成的 \(Z(\lambda)\) 的子模。则 \(x \neq 0\) （命题 6）。另一方面，\(x \in Z(\lambda)^{\lambda-(m+1)\alpha}\) 。对 \(\beta \in B\) 且 \(\beta \neq \alpha\) ，有 \([\mathfrak{g}^{-\alpha}, \mathfrak{g}^{\beta}] = 0\) 且 \(g^{\beta}.e = 0\) ，故 \(g^{\beta}.x = 0\) 。最后，由于 \([X_{\alpha}, X_{-\alpha}] = -H_{\alpha}\) ，我们有

\[
[ X _ {\alpha}, X _ {- \alpha} ^ {m + 1} ] = (m + 1) X _ {- \alpha} ^ {m} (- H _ {\alpha} + m)
\]

（§1 第 1 节 引理 1 (ii)），故

\[
X _ {\alpha}. x = X _ {\alpha} X _ {- \alpha} ^ {m + 1}. e = [ X _ {\alpha}, X _ {- \alpha} ^ {m + 1} ]. e = (m + 1) X _ {- \alpha} ^ {m} (m e - \lambda (H _ {\alpha}) e) = 0.
\]

于是，\(x\) 是权为 \(\lambda - (m + 1)\alpha\) 的本原元。鉴于命题 6，\(\mathfrak{g}\) -模 \(V\) 同构于 \(Z(-\alpha + \lambda - m\alpha) = Z(-\alpha + s_{\alpha}\lambda)\) 。

推论 2. 设 \(\rho = \frac{1}{2}\sum_{\alpha \in \mathrm{R}_{+}}\alpha\) ，并设 \(\lambda, \mu \in \mathfrak{h}^*\) 。假设 \(\lambda + \rho\) 是 \(\mathrm{R}\) 中的一个支配权，且存在 \(w \in \mathrm{W}\) 使 \(\mu + \rho = w(\lambda + \rho)\) 成立。则 \(\mathrm{Z}(\mu)\) 同构于 \(\mathrm{Z}(\lambda)\) 的一个子模。

当 w = 1 时断言是明显的。假设当 w 的长度 \textless{} q 时断言已经成立。若 w 的长度为 q，则存在 \(\alpha \in B\) 使 \(w = s_{\alpha}w'^{-1}\) 成立，其中 \(l(w') = q - 1\) 。我们有 \(w'(\alpha) \in \mathrm{R}_{+}\) （第六章 §1 第 6 节 命题 17 的推论 2），因而 \(w'^{-1}(\lambda + \rho)(H_{\alpha}) = (\lambda + \rho)(H_{w'\alpha})\) 是一个 \(\geq 0\) 的整数。置

\[
\mu^ {\prime} = w ^ {- 1} (\lambda + \rho) - \rho .
\]

由归纳假设，\(Z(\mu')\) 同构于 \(Z(\lambda)\) 的一个子模。另一方面，由第六章 §1 第 10 节 命题 29 (ii)，

\[
- \alpha + s _ {\alpha} \mu^ {\prime} = - \alpha + s _ {\alpha} w ^ {- 1} (\lambda + \rho) - s _ {\alpha} \rho = w (\lambda + \rho) - \rho = \mu .
\]

此外，\(\rho(H_{\alpha}) = 1\) （第六章 §1 命题 29 (iii)），故 \(\mu'(H_{\alpha}) + 1 \in \mathbf{N}\) 。于是由推论 1 知，\(Z(\mu)\) 同构于 \(Z(\mu')\) 的一个子模，从而也同构于 \(Z(\lambda)\) 的一个子模。

